\scanpage{69}
il semble\inserted{rait} que la structure \unsure{$S$}
\[
 \left\{\begin{aligned}F_M&\subset\mathfrak P_3(S_M)\\ F_M&\xrightarrow{\tau_F}\unsure{\pi|\Pi}_M\end{aligned}\right.
\]
soit nettement plus fine que la structure $A_M\subset\mathfrak P_2(S_M),\ A_M\xrightarrow{\tau_A}\pi_M$ qui s'en déduit, via (177) (168) :
\deleted{Une application de $\bar M$ \unsure{supposons}}
$\Lambda$ un corps $k$ de car. 2, \deleted{\unsure{$\ne$}} corps premier $\mathbb F_2$, alors une \deleted{bijection}\inserted{permutation} de \deleted{\unsure{$M$}}\inserted{$M^*=M\smallsetminus\{0\}$ (où $M=k^2$)} peut fort bien transformer couples de pts \deleted{\unsure{qui sont vus en|qui sont une base}} \deleted{\unsure{par}} \deleted{alignés}\inserted{colinéaires} en couples \deleted{alignés}\inserted{colinéaires} --- donc aussi bases en bases --- et, je suppose, de façon compatible avec les jauges i.e. en transformant \deleted{\unsure{base}} deux bases de \unsure{$=$|même} jauge en \deleted{deux} \inserted{des} bases de $=$ jauge, sans pour autant que \unsure{respecte|respecter} la relation $x+y+z=0$, pour \deleted{\unsure{$x,y,z$}}, (\inserted{avec} $x\wedge y=y\wedge z=z\wedge x\in(\det M)^*$) \ldots.
% slanted
\margin{\uline{à élucider} éventuellement, en utilisant la structure d'addition non partout définie de $\pi_M=(\det M)^*$, \deleted{\unsure{en}} et exigeant qu'elle soit respectée par les \deleted{\unsure{jauges}} \unsure{autom|permutations}, on arrive à déduire parfois $F_M$ de $A_M,\ \tau_A:A_M\to\pi_M$ \ldots}

À vrai dire, la structure que nous étudions en car. 2, se réduisant \inserted{alors} à l'ens. des sommets $\widetilde S_M=M^{**}$ (ens. des éléments de $M$ faisant partie d'une base), l'ens. $\widetilde F_M\unsure{\subset|\not\subset|\notin}\mathfrak P_3(S)$, formé de parties \inserted{$\{x,y,z\}$} de $\widetilde S_M$ satisfaisant $x+y+z=0$, $x\wedge y\ (=y\wedge z=z\wedge x)\in(\det M)^*$, et de l'application
\[
 \widetilde\tau:\widetilde F_M\longrightarrow(\det M)^*=\widetilde\pi\quad
 \{x,y,z\}\longmapsto\begin{aligned}[t]x\wedge y&=y\wedge z\\&=z\wedge x\end{aligned}\tag*{(178)}
\]
% archive sheet number: 68
